\section{Lecture 19: Regularization in Deep Learning}

\subsection{Motivation}

Deep neural networks are highly expressive and often overparameterized.

\medskip

\noindent
As a result:
\begin{itemize}
    \item they can fit training data extremely well,
    \item but may overfit and generalize poorly.
\end{itemize}

\medskip

\noindent
\textbf{Key question.}  
How can we control model complexity and improve generalization in deep learning?

\medskip

\noindent
\textbf{Guiding idea.}  
Regularization techniques constrain or perturb learning to favor simpler and more robust solutions.

\subsection{Dropout and Noise Injection}

\textbf{Dropout.}

During training, randomly set activations to zero:

\[
h_i \mapsto h_i \cdot \xi_i, \quad \xi_i \sim \text{Bernoulli}(p).
\]

\medskip

\noindent
\textbf{Effect.}
\begin{itemize}
    \item Prevents co-adaptation of neurons
    \item Encourages redundancy and robustness
\end{itemize}

\medskip

\noindent
\textbf{Inference.}
\begin{itemize}
    \item Use full network with scaled activations
\end{itemize}

\medskip

\noindent
\textbf{Interpretation.}
\begin{itemize}
    \item Approximate model averaging over subnetworks
\end{itemize}

\medskip

\noindent
\textbf{Noise injection.}

More generally, we can inject noise into:

\begin{itemize}
    \item inputs,
    \item activations,
    \item gradients.
\end{itemize}

\medskip

\noindent
\textbf{Insight.}  
Noise prevents over-reliance on specific features and improves robustness.

\subsection{Weight Decay}

Weight decay adds an $\ell_2$ penalty:

\[
\min_\theta \; L(\theta) + \lambda \|\theta\|^2.
\]

\medskip

\noindent
\textbf{Interpretation.}
\begin{itemize}
    \item Penalizes large weights
    \item Encourages simpler models
\end{itemize}

\medskip

\noindent
\textbf{Connection to optimization.}
\begin{itemize}
    \item Equivalent to shrinking parameters during updates
\end{itemize}

\medskip

\noindent
\textbf{Geometric view.}
\begin{itemize}
    \item Constrains solution to lie in a bounded region
\end{itemize}

\medskip

\noindent
\textbf{Insight.}  
Weight decay controls model capacity through parameter magnitude.

\subsection{Data Augmentation}

\textbf{Idea.}  
Generate additional training data by applying transformations.

\medskip

\noindent
\textbf{Examples.}
\begin{itemize}
    \item Image transformations:
    \begin{itemize}
        \item rotation, cropping, flipping
    \end{itemize}

    \item Text transformations:
    \begin{itemize}
        \item synonym replacement, paraphrasing
    \end{itemize}
\end{itemize}

\medskip

\noindent
\textbf{Interpretation.}
\begin{itemize}
    \item Encodes invariances in the data
\end{itemize}

\medskip

\noindent
\textbf{Effect.}
\begin{itemize}
    \item Expands effective dataset
    \item Improves generalization
\end{itemize}

\medskip

\noindent
\textbf{Insight.}  
Data augmentation injects prior knowledge into learning.

\subsection{Implicit Bias of Optimization}

Even without explicit regularization, optimization itself induces bias.

\medskip

\noindent
\textbf{Examples.}

\begin{itemize}
    \item Gradient descent in linear models:
    \begin{itemize}
        \item finds minimum-norm solution
    \end{itemize}

    \item SGD in deep networks:
    \begin{itemize}
        \item prefers flatter minima
    \end{itemize}
\end{itemize}

\medskip

\noindent
\textbf{Flat vs sharp minima.}

\begin{itemize}
    \item Flat minima:
    \begin{itemize}
        \item robust to perturbations
        \item better generalization
    \end{itemize}

    \item Sharp minima:
    \begin{itemize}
        \item sensitive to noise
        \item poorer generalization
    \end{itemize}
\end{itemize}

\medskip

\noindent
\textbf{Role of SGD.}
\begin{itemize}
    \item Noise in updates biases toward flat regions
\end{itemize}

\medskip

\noindent
\textbf{Insight.}  
Optimization dynamics act as an implicit regularizer.

\subsection{A Unifying Perspective}

Regularization in deep learning arises from multiple sources:

\begin{itemize}
    \item \textbf{Explicit:}
    \begin{itemize}
        \item weight decay, dropout
    \end{itemize}

    \item \textbf{Data-driven:}
    \begin{itemize}
        \item augmentation
    \end{itemize}

    \item \textbf{Implicit:}
    \begin{itemize}
        \item optimization dynamics
    \end{itemize}
\end{itemize}

\medskip

\noindent
\textbf{Core idea.}  
Generalization is controlled not only by model size, but by how the model is trained.

\medskip

\noindent
\textbf{Key insight.}  
Modern deep learning relies heavily on implicit and data-driven regularization.

\subsection{Outlook}

In this lecture, we studied:

\begin{itemize}
    \item dropout and noise injection,
    \item weight decay,
    \item data augmentation,
    \item implicit bias of optimization.
\end{itemize}

\medskip

\noindent
These techniques improve generalization in deep networks.

\medskip

\noindent
In the next lecture, we study residual networks and skip connections, which enable training of very deep architectures.

\medskip

\noindent
\textbf{Guiding principle.}  
Effective deep learning balances expressivity, stability, and generalization.